\section{Introduction}
\label{sec:intro}
Innate immune cells keep a memory of earlier stimuli. Monocytes and macrophages exposed to
$\beta$-glucan or BCG respond more strongly to an unrelated second stimulus days later
(trained immunity), while cells exposed to LPS respond less (tolerance); both states are
carried by epigenetic and metabolic reprogramming rather than by antigen-specific receptors
\citep{netea2016,saeed2014,novakovic2016}. These states are measured today by restimulating
cells and assaying cytokines, transcripts or chromatin, which destroys the sample and gives
one time point. If memory states left a signature in how living cells look and move, label-free
video could read them out continuously and non-destructively.

Self-supervised video models make that question testable without labels: a joint-embedding
predictive architecture (JEPA) learns representations by predicting the latent content of
masked or future parts of a clip \citep{vjepa,vjepa2}. But testing whether such a model
captures immune state requires paired data, images or video of cells whose molecular or
functional state was measured, and for trained innate cells none is public. This paper asks
how far one can get with public data only. We train TiJepa, a video JEPA, on
\val{pre.t1.clips} unlabelled clips of migrating human neutrophils, and test its frozen
representations, against strong image and video models, on every public dataset we could
find that pairs label-free images, video or image-derived dynamics of innate immune cells with a
molecular, functional or memory readout.

The contribution is an audit rather than a new state of the art, and it has three parts.
\begin{enumerate}
\item A \emph{modality coverage map} (Figure~\ref{fig:coverage}, Table~\ref{tab:coverage}):
for each modality family (immunophenotype, metabolism, signalling, transcriptomics,
epigenomics, cytokines, integrated omics), the best public pairing with imaging, the level
at which it is paired (same cell, same condition, different studies), what we tested and
the result, including the families for which no public pairing exists.
\item A controlled evaluation of what a video JEPA learns from innate-immune video: a
matched-step spatial-only ablation, a frozen general-purpose video model (V-JEPA~2), image
models, random-initialisation and frame-shuffle controls, a label-free validation set used
for every model choice, controls that separate forecasting from positional structure in the
latent space, and a diagnosis of the partial dimensional collapse we observed.
\item From the molecular data that exist without images, a replicated observation that
constrains what an imaging study of trained cells should measure, and the design of that
study.
\end{enumerate}
Most of our answers are negative or null, and we report them as such. No result in this
paper shows trained immunity in images: no public image of trained cells exists, and the
memory-related analyses use the closest public proxies (acute priming and tolerance,
polarisation, restimulation transcriptomes).
